\chapter{LITERATURE REVIEW}

\section{Computational Approaches in Hydrocarbon Prediction}

The interpretation of well log and mud log data has historically depended on the expertise and judgment of individual petrophysicists. As reservoirs grow more complex and the volume of acquired data continues to expand, the limitations of purely manual interpretation have become increasingly apparent. Human analysis is inherently time-consuming, difficult to scale, and subject to variability between interpreters. Computational approaches address these limitations directly: they offer systematic, reproducible pipelines capable of processing large multi-well datasets with consistent logic and quantifiable uncertainty. Two broad paradigms have emerged as dominant methodological frameworks in this space --- deterministic algorithmic approaches rooted in established physical laws, and data-driven approaches based on machine learning and deep learning architectures. Both have been demonstrated in the literature to provide meaningful uplift over manual interpretation alone, and an understanding of their respective strengths and constraints is essential for contextualizing the methodology of this study.

\subsection{Deterministic and Algorithmic Approaches}

Deterministic approaches to formation evaluation derive their outputs through the direct application of explicit mathematical models to measured log data. The intellectual foundation of this paradigm can be traced to the seminal work of \citet{archie1942}, who established the empirical relationship between a formation's measured electrical resistivity, its porosity, and the saturation of its pore fluids. This equation, now universally known as Archie's Law, provided the first rigorous quantitative framework for inferring subsurface fluid content from borehole measurements, and it remains a cornerstone of formation evaluation to this day. Subsequent decades brought additional deterministic models addressing shaly sand corrections (the Simandoux equation), volumetric mineral decomposition, and the interpretation of nuclear and acoustic logs --- all of which operate on the same fundamental principle: given a set of calibrated input measurements and known physical relationships, the output can be computed directly and reproduced exactly \citep{ref9, ref11}.

A defining characteristic of deterministic methods is their full transparency. Every intermediate result can be audited, the sensitivity of the output to each input parameter can be quantified, and any experienced petrophysicist can trace the chain of reasoning from raw log reading to final reservoir classification. This auditability is not merely an academic advantage; it is often a regulatory and commercial requirement. The \citet{ref3} standards for petroleum resource estimation, for instance, explicitly distinguish deterministic from probabilistic assessment procedures and require that the methodology underlying any reserves declaration be clearly documented and defensible. A code-based deterministic system --- one in which Archie's Law, shale volume corrections, and pay-zone criteria are encoded as explicit algorithmic rules --- inherently satisfies this requirement \citep{ref11}.

The practical implementation of deterministic petrophysical logic in software has a long history. Early systems encoded expert knowledge as lookup tables and threshold-based decision trees. More sophisticated modern implementations use layered conditional logic that integrates multiple log curves simultaneously --- for instance, cross-referencing gamma ray for lithology discrimination, resistivity for fluid identification, and neutron-density crossplots for porosity calibration --- to classify each depth interval. The transparency of such rule-based systems means that domain experts can inspect, challenge, and refine the underlying logic iteratively, which is particularly valuable in frontier basins where geological priors are poorly constrained. The primary limitation of purely deterministic approaches is their rigidity: the accuracy of the output is bounded by the validity of the underlying physical model, and classical equations like Archie's Law carry known assumptions (clean formations, water-wet pores, homogeneous lithology) that are frequently violated in complex reservoirs \citep{ref9, ref11}.

\subsection{Machine Learning and Data-Driven Approaches}

Parallel to deterministic methods, the application of machine learning (ML) and deep learning (DL) to formation evaluation and hydrocarbon prediction has grown substantially since its early introduction to petroleum engineering. The conceptual groundwork for applying artificial intelligence techniques to oilfield problems was laid by \citet{ref1}, who demonstrated that artificial neural networks could approximate complex, non-linear mappings between well log inputs and reservoir properties that resist explicit physical formulation. That early work catalysed a sustained research trajectory that has since expanded to include Random Forests, Support Vector Machines, gradient-boosted tree ensembles, and convolutional and recurrent neural network architectures.

The primary appeal of ML approaches is their ability to learn structure directly from data without requiring the analyst to specify the functional form of the relationship between input and output. In heterogeneous reservoirs where the porosity--permeability relationship is non-stationary, where lithological facies change rapidly along the wellbore, or where the applicability of classical deterministic equations is uncertain, ML models can often achieve higher predictive accuracy than their physics-based counterparts. This has been demonstrated across a range of formation evaluation tasks. \citet{ref5} showed that ML models trained on wireline log suites can predict in-situ geomechanical properties with strong accuracy, highlighting the utility of data-driven methods in applications where direct measurement is prohibitively expensive. The task of reconstructing missing or degraded log intervals --- a common practical problem in older wells or in wells where certain tools were not run --- has similarly been addressed through DL architectures. \citet{ref6} introduced a flexible graph-neural-network approach, FlexLogNet, capable of predicting any arbitrary missing log from the remaining available logs within the same well, adaptively adjusting its input configuration to whatever data is present. Well-log completion methods of this type have become increasingly important as operators attempt to extract maximum value from legacy datasets containing incomplete log suites \citep{ref4, ref6}.

Despite these demonstrated capabilities, ML and DL approaches carry inherent limitations that are particularly consequential in a high-stakes petroleum context. The most widely cited of these is the interpretability problem: because complex neural networks and ensemble models derive their predictions through high-dimensional, non-linear transformations that are not straightforwardly human-readable, it can be difficult or impossible to explain why a specific prediction was generated. This opacity creates friction in professional and regulatory settings, where the basis for any reservoir characterization must be auditable \citep{ref8}. A second major constraint is data dependency: supervised ML models require sufficient labelled training examples to generalise reliably to unseen wells. In frontier exploration areas, where few wells have been drilled and production data is scarce, the training data volumes required to fit complex models may simply not be available, limiting their practical applicability \citep{ref1}. A third concern is distribution shift: models trained on data from one geological province may perform poorly when applied to a different basin with distinct lithological characteristics, even if the surface log signatures appear superficially similar.